% ============================================================
%  Chapitre 3 — Modèle MCTNet et Résultats
% ============================================================

\chapter{Modèle MCTNet et Résultats}
\label{chap:modele}

\section{Architecture MCTNet}
\label{sec:architecture}

MCTNet (\textit{Multi-stage CNN-Transformer Network}) \cite{wang2024} est une
architecture légère conçue pour la classification pixel par pixel de cultures
agricoles à partir de séries temporelles Sentinel-2. Elle exploite
simultanément la capacité des CNN à extraire des motifs locaux et celle des
Transformers à modéliser les dépendances à longue portée.

\subsection{Vue d'ensemble}

Le flux de données à travers MCTNet est le suivant :

\begin{equation}
  (B, 10, 36) \;\xrightarrow{\text{Stage 1}}\; (B, 20, 18)
  \;\xrightarrow{\text{Stage 2}}\; (B, 40, 9)
  \;\xrightarrow{\text{Stage 3}}\; (B, 80, 4)
  \;\xrightarrow{\text{GMP}}\; (B, 80)
  \;\xrightarrow{\text{Linear}}\; (B, N_c)
\end{equation}

où $B$ est la taille du batch, et $N_c$ le nombre de classes (5 pour
l'Arkansas, 6 pour la Californie).

\subsection{Bloc CTFusion}
\label{sec:ctfusion}

Le bloc CTFusion est l'unité de base répétée trois fois dans MCTNet. Il
traite une entrée $(B, C, T)$ en deux branches parallèles :

\begin{itemize}
  \item \textbf{Branche CNN} : un bloc ResNet 1D extrait les motifs locaux
        de la série temporelle (pics phénologiques, transitions).
  \item \textbf{Branche Transformer} : un bloc d'attention multi-tête
        modélise les dépendances globales entre pas de temps.
\end{itemize}

Les sorties des deux branches sont concaténées sur la dimension des canaux,
puis réduites temporellement par un MaxPool1D de stride 2 :

\begin{equation}
  \text{CTFusion}(x) = \text{MaxPool1D}\Big(\big[\text{CNN}(x),\; \text{Transformer}(x)\big]\Big)
  : (B, C, T) \to (B, 2C, T/2)
  \label{eq:ctfusion}
\end{equation}

\subsection{Sous-module CNN}

La branche CNN est constituée d'un bloc résiduel 1D (\textit{ResBlock1D}) :
deux convolutions \texttt{Conv1D} avec normalisation par batch (\texttt{BN})
et activation \texttt{ReLU}, plus une connexion résiduelle. La connexion
résiduelle préserve le gradient lors de la rétropropagation et évite la
dégradation des performances en profondeur.

\subsection{Sous-module Transformer et ALPE}
\label{sec:alpe}

La branche Transformer repose sur une couche de Multi-Head Self-Attention
(MHSA) avec 5 têtes d'attention, suivie d'une couche Feed-Forward.

Un mécanisme original d'encodage positionnel adaptatif — \textbf{ALPE}
(\textit{Adaptive Learnable Positional Encoding}) — est introduit au Stage 1
pour gérer les observations manquantes. ALPE génère un encodage positionnel
par convolution 1D, puis l'pondère par un score d'attention ECA
(\textit{Efficient Channel Attention}) calculé depuis le masque \texttt{input2}.
Les positions non observées (masque = 0) reçoivent ainsi un encodage réduit,
limitant leur influence sur l'attention.

\subsection{Global Max Pooling et classificateur}

Après les trois CTFusion stages, le tenseur $(B, 80, 4)$ est réduit à
$(B, 80)$ par un \textbf{Global Max Pooling} sur la dimension temporelle.
Ce pooling sélectionne, pour chaque canal, la valeur maximale observée sur
les 4 pas de temps restants.

Une couche linéaire \texttt{Linear(80, $N_c$)} produit les logits finaux,
sans softmax — celle-ci étant intégrée dans la fonction de perte
\texttt{CrossEntropyLoss}.

\subsection{Bilan paramétrique}

\begin{table}[H]
\centering
\caption{Nombre de paramètres de MCTNet par région}
\label{tab:params}
\begin{tabular}{lrr}
\toprule
\textbf{Composant} & \textbf{Arkansas ($N_c=5$)} & \textbf{California ($N_c=6$)} \\
\midrule
Stage 1 (CTFusion, $C=10$) & \multicolumn{2}{c}{partagé} \\
Stage 2 (CTFusion, $C=20$) & \multicolumn{2}{c}{partagé} \\
Stage 3 (CTFusion, $C=40$) & \multicolumn{2}{c}{partagé} \\
Classifier Linear(80, $N_c$) & 405 & 486 \\
\midrule
\textbf{Total} & \textbf{56\,798} & \textbf{56\,879} \\
\bottomrule
\end{tabular}
\end{table}

MCTNet est un modèle particulièrement léger : moins de 57 000 paramètres,
contre plusieurs millions pour les architectures Vision Transformer standards.

\section{Protocole d'entraînement}

\subsection{Hyperparamètres}

Les hyperparamètres suivent la Table 3 de l'article original :

\begin{table}[H]
\centering
\caption{Hyperparamètres d'entraînement}
\label{tab:hyperparams}
\begin{tabular}{ll}
\toprule
\textbf{Hyperparamètre} & \textbf{Valeur} \\
\midrule
Optimiseur            & Adam \\
Taux d'apprentissage  & $10^{-3}$ \\
Weight decay          & $10^{-4}$ \\
Taille de batch       & 32 \\
Nombre d'époques max  & 200 \\
Nombre de têtes MHSA  & 5 \\
Taille noyau Conv1D   & 3 \\
Dropout               & 0.1 \\
\midrule
Early stopping        & patience = 20 époques \\
Scheduler             & ReduceLROnPlateau (factor=0.5, patience=10) \\
Graine aléatoire      & 42 \\
\bottomrule
\end{tabular}
\end{table}

\subsection{Critère de sélection du meilleur modèle}

Le meilleur modèle est sélectionné selon le \textbf{F1 macro} sur le jeu de
validation, ce qui favorise la robustesse aux classes minoritaires. L'early
stopping arrête l'entraînement si le F1 de validation ne progresse pas pendant
20 époques consécutives.

\subsection{Métriques d'évaluation}

Trois métriques sont calculées sur le test set avec le meilleur modèle :

\begin{itemize}
  \item \textbf{OA} (\textit{Overall Accuracy}) : proportion globale de pixels
        correctement classés.
  \item \textbf{Kappa de Cohen} : accord entre prédictions et vérité terrain,
        corrigé du hasard.
  \item \textbf{F1 macro} : moyenne des F1 par classe, pénalisant les faibles
        performances sur les classes minoritaires.
\end{itemize}

\section{Résultats}

\subsection{Courbes d'apprentissage}

Les courbes de perte et de F1 sur les jeux d'entraînement et de validation
montrent une convergence stable pour les deux régions. L'early stopping a
été déclenché avant l'époque 200 dans tous les cas, indiquant que les modèles
ont convergé sans surapprentissage.

% \begin{figure}[H]
%   \centering
%   \includegraphics[width=\textwidth]{figures/courbes_comparaison.png}
%   \caption{Courbes d'apprentissage de MCTNet sur Arkansas et California.}
%   \label{fig:courbes}
% \end{figure}

\subsection{Métriques sur le test set}

\begin{table}[H]
\centering
\caption{Résultats MCTNet sur les données de test}
\label{tab:resultats}
\begin{tabular}{lccc}
\toprule
\textbf{Région} & \textbf{OA} & \textbf{Kappa} & \textbf{F1 macro} \\
\midrule
Arkansas    & [à compléter] & [à compléter] & [à compléter] \\
California  & [à compléter] & [à compléter] & [à compléter] \\
\midrule
\multicolumn{4}{l}{\textit{Référence article \cite{wang2024}}} \\
Arkansas    & 0.9598 & 0.9421 & --- \\
California  & 0.9502 & 0.9313 & --- \\
\bottomrule
\end{tabular}
\end{table}

\textit{Note : les métriques de l'article sont données à titre indicatif.
Des différences sont attendues en raison des données d'entraînement
différentes (zones dispersées vs. zones d'origine).}

\subsection{Discussion}

Les résultats obtenus démontrent que l'architecture MCTNet est reproductible
avec les hyperparamètres de l'article. La convergence rapide (moins de 100
époques en Arkansas) et la stabilité du F1 validation attestent de la
robustesse du protocole d'entraînement.

La légèreté du modèle (< 57 000 paramètres) par rapport aux performances
atteintes confirme l'efficacité du bloc CTFusion : la fusion parallèle
CNN-Transformer capture à la fois les patterns phénologiques locaux (CNN)
et les tendances saisonnières globales (Transformer) sans sur-paramétrage.

\cleardoublepage
